\scanpage{4}
\mspage{1 bis}
\uline{Lemme} $A,B$ schémas abéliens sur $S$ loc. noeth, $u_\ell:T_\ell(A)\unsure{-|\longrightarrow}T_\ell(B)$. Alors $\exists$ \uline{sous-schéma fermé} $T$ de $S$ tel que, pour tout $S'/S$, on \inserted{\unsure{ceci|ait ceci|ait} :} $u_{\ell S'}$ provient de $u:A_{S'}\unsure{-|\longrightarrow}B_{S'}$ \unsure{sss|ssi} $S'$ se factorise par $T$.

Cf Inventions 1.2 (\unsure{ou|où|une} démonstration \unsure{cristalline|conditionnelle}).

\deleted{\uline{Cor} Supposons que $\forall s\in\unsure{S|\operatorname{Ass}\mathcal O_S}$, \unsure{l'une des|l'image de} hyp. suiv}

\uline{Complément} De plus, \deleted{\unsure{de la conclusion}} pour tt $t\in T$ tel que $\mathrm{car.}\,k(t)=\ell$, $T$ majore un voisinage ouvert immédiat de $t$.

\uline{Corollaire} Supposons que l'ouvert $U\subset S$ soit tel que

a) tte comp. conn. de $S$ rencontre $U$

b) $t\in S-U\Longrightarrow\mathrm{car.}\,k(t)=\ell$.

c) $u_\ell|U$ provient de $u_U:A/U\unsure{-|\longrightarrow}B/U$

Alors \deleted{\unsure{toute}} $u_\ell$ provient de $u:A\longrightarrow B$

\deleted{En effet, on peut supposer $S$ connexe, \inserted{\unsure{donc} a) $\Longleftrightarrow U\neq\varnothing$.} Alors \deleted{$\unsure{T\text{ est }S}$} \unsure{tout pt de} \deleted{\unsure{l'une des} composantes irréductibles} [distinguer $t\in\deleted{\unsure{T}}U$ et $t\in T\cap(S-U)$ \inserted{grâce à b)}], \deleted{\unsure{de} $S$. Si on avait \unsure{disons} $T\neq S$ on aurait une comp. $S_i$ de $S$ telle} Comme $T\neq\varnothing$, il suffit de prouver que tte comp. irréd. $S_i$ de $S$ qui rencontre $T$ est $\subset T$. Or si $S_i\cap U=\varnothing$ \unsure{clair|clairement}, \inserted{car $T\supset\overline U$} \unsure{sinon} aussi grâce au complément \inserted{et à b)}. Donc $T\supset S_{\mathrm{red}}$. \deleted{\unsure{Ainsi}} Mais $T$ est \deleted{\unsure{indu|fermé}} ouvert} $T$ est ouvert \unsure{évident|induit} sur $U$ par hyp. \inserted{\deleted{et} c)}, et en les pts de $(S-U)\cap T$ grâce au complément et à b). Donc $\deleted{U}\unsure{T|\overline U|\overline T}$ est un \inserted{sous-schéma} ouvert, et comme il est fermé, c'est une réunion de comp. connexes. Par a) on a donc $T=S$, cqfd

% archive sheet number: 3
